\section{Conclusion}\label{sec:concl}
This study investigated whether introducing bounded-authority agent specialization into an adaptive technical-interview system changes decision trajectories, and whether observed changes stem from specialization itself or from the coordination and arbitration layer. Across 160 deterministic replay sessions, the multi-agent configuration altered difficulty trajectories relative to the frozen single-policy baseline, exhibiting fewer level changes and reversals, although leave-one-cluster-out analysis showed that the magnitude was sensitive to the eight replay clusters, alongside $+0.59$ proactive probe turns per session that arose only in replay cluster $k=1$. In the tested replay, the ablations associate this behavioral change with the Strategy proposal and confidence-gated arbitration path. In contrast, matched open-loop comparison between the specialized two-agent architecture and a unified controller produced identical decisions across all 1,280 turns ($\mathbb{P}(\text{agree}) = 1.0$), indicating that role decomposition did not inherently alter decision behavior under equivalent heuristic rules and observations.

The evaluation does not establish that reduced trajectory volatility reflects superior adaptation, nor does it demonstrate candidate learning gains, pedagogical efficacy, or general equivalence for stochastic LLM agents. Instead, the evidence supports an architectural distinction between specialization as a mechanism for structuring responsibility and private state, and coordination and arbitration as mechanisms that govern proposal admission into committed decisions. Bounded proposal contracts, confidence gating, and a centralized commit hub in \texttt{InterviewOrchestrator} allow multi-agent interaction to be structured and audited without decentralized state mutation. Future work should evaluate whether these boundaries persist across diverse learned policy checkpoints, broader candidate populations, closed-loop human interaction, and generative multi-agent architectures.
